\documentclass[11pt,a4paper]{article}

\usepackage[utf8]{inputenc}
\usepackage[T1]{fontenc}
%\usepackage[spanish]{babel}
\usepackage{amsmath,amssymb,bm,amsthm,mathtools}
\usepackage{geometry}
\usepackage{booktabs}
\usepackage{array,tabularx}
\usepackage{float}
\usepackage{enumitem}
\usepackage{xcolor}
\usepackage{tcolorbox}
\usepackage{hyperref}
\usepackage{caption}
\usepackage{tikz}
\usetikzlibrary{arrows.meta,positioning,calc}

\geometry{margin=2.5cm}

\definecolor{cblue}{RGB}{41,128,185}
\definecolor{cgreen}{RGB}{39,174,96}
\definecolor{cred}{RGB}{231,76,60}
\definecolor{corange}{RGB}{243,156,18}
\definecolor{cpurple}{RGB}{142,68,173}
\definecolor{cdark}{RGB}{44,62,80}

\tcbset{
  radio/.style={colback=corange!5,colframe=corange!60,fonttitle=\bfseries,arc=2mm,boxrule=0.5pt},
  math/.style={colback=cblue!5,colframe=cblue!60,fonttitle=\bfseries,arc=2mm,boxrule=0.5pt},
  result/.style={colback=cgreen!5,colframe=cgreen!60,fonttitle=\bfseries,arc=2mm,boxrule=0.5pt},
  warning/.style={colback=cred!5,colframe=cred!60,fonttitle=\bfseries,arc=2mm,boxrule=0.5pt},
  key/.style={colback=cpurple!5,colframe=cpurple!60,fonttitle=\bfseries,arc=2mm,boxrule=0.5pt},
}

\newtheorem{definition}{Definici\'on}[section]
\newtheorem{theorem}{Teorema}[section]
\newtheorem{proposition}{Proposici\'on}[section]

\newcommand{\pw}{P(\text{win})}
\newcommand{\E}{\mathbb{E}}
\newcommand{\Prob}{\mathbb{P}}

\begin{document}

\begin{center}
{\LARGE\bfseries Framework de Dise\~no de Estrategias}\\[4pt]
{\LARGE\bfseries Basado en Bimodalidad y Teor\'ia de la Informaci\'on}\\[16pt]
{\large Javier Epeloa\quad$\cdot$\quad Mapplics}\\[8pt]
{\normalsize Abril 2026}\\[4pt]
{\small\itshape Metodolog\'ia completa: de la se\~nal candidata al sistema de salida \'optimo\\
con un solo par\'ametro libre, sin backtesting, sin optimizaci\'on}
\end{center}

\vspace{12pt}

\begin{abstract}
Este documento describe un framework de cinco pasos para dise\~nar
sistemas de trading en cualquier mercado con microestructura observable.
El framework reduce un problema de 14 par\'ametros calibrados a un sistema
con un solo par\'ametro libre, una tabla emp\'irica, y un test estad\'istico
de falsificaci\'on. A lo largo del documento, cada concepto matem\'atico se
acompa\~na de una analog\'ia con un receptor de radio AM que sintoniza
estaciones, decodifica mensajes, y decide cu\'ando dejar de escuchar.
\end{abstract}

\tableofcontents
\newpage

%=================================================================
\section{La Analog\'ia Central: El Trade como Canal de Radio}
\label{sec:analogia}
%=================================================================

Antes de entrar en la matem\'atica, necesitamos una imagen mental
que nos acompa\~ne durante todo el documento.

\begin{tcolorbox}[radio, title=La radio]
Imagin\'a que ten\'es una radio vieja con dial anal\'ogico. El dial tiene
infinitas posiciones --- la mayor\'ia son ruido blanco, est\'atica pura.
Pero en algunas frecuencias hay estaciones que transmiten una se\~nal.

El problema es que no sab\'es d\'onde est\'an las estaciones. Y cuando
encontr\'as una, no sab\'es qu\'e mensaje est\'a transmitiendo hasta que
escuch\'as un rato.
\end{tcolorbox}

Cada componente del framework tiene su equivalente en la radio:

\begin{table}[H]
\centering
\renewcommand{\arraystretch}{1.4}
\begin{tabularx}{\textwidth}{lX}
\toprule
\textbf{Radio} & \textbf{Trading} \\
\midrule
Girar el dial & Probar una se\~nal de entrada candidata \\
?`Hay estaci\'on? & Test de bimodalidad ($\Delta$BIC $> 10$) \\
Se\~nal vs.\ ruido & Winners vs.\ losers en la distribuci\'on de MFE \\
Escuchar el mensaje & Observar la trayectoria del precio tick a tick \\
Decodificar W o L & Calcular $\pw$ desde la superficie \\
Dejar de escuchar & Cerrar el trade (EXIT o TRAIL) \\
La estaci\'on se debilita & El edge se erosiona ($\Delta$BIC cae) \\
\bottomrule
\end{tabularx}
\end{table}


%=================================================================
\section{Paso 1: Girar el Dial --- Se\~nal Candidata en Dry Mode}
\label{sec:paso1}
%=================================================================

\begin{tcolorbox}[radio, title=La radio: girar el dial]
Mov\'es el dial a una frecuencia aleatoria. No sab\'es si hay una
estaci\'on ah\'i. Solo escuch\'as --- sin tocar ning\'un otro bot\'on,
sin ajustar el volumen, sin nada. Escuch\'as un rato fijo (por ejemplo,
4 minutos) y anot\'as todo lo que o\'is. Despu\'es mov\'es el dial
a otra frecuencia y repet\'is.
\end{tcolorbox}

\subsection{Qu\'e hacer}

Tom\'as cualquier se\~nal candidata --- puede ser RSI $> 80$, un cruce
de medias m\'oviles, una se\~nal de funding rate, un modelo de ML, lo que
sea. Abr\'is trades virtuales (dry mode) cada vez que la se\~nal
se activa. Sin stop loss, sin trailing, sin ninguna gesti\'on. El trade
abre, corre un horizonte fijo (por ejemplo, 4 horas), y cierra.

Durante esas 4 horas, logue\'as el \textbf{path data}: cada 10--30
segundos, registr\'as MFE (M\'axima Excursi\'on Favorable), MAE (M\'axima
Excursi\'on Adversa), y PnL actual.

\subsection{Por qu\'e horizonte fijo}

El horizonte fijo es cr\'itico porque elimina la circularidad. Si us\'as
stop loss o trailing, el resultado del trade depende del mecanismo de
salida, no solo de la se\~nal. No pod\'es saber si la se\~nal es buena o
si la gesti\'on es buena --- est\'an confundidas.

Con horizonte fijo, la definici\'on de winner es pura:
\begin{equation}
\text{Winner} = \mathbb{1}[\text{PnL}(T_{\text{horizonte}}) > 0]
\end{equation}

No hay par\'ametros. No hay decisiones. Solo datos limpios.

\subsection{Justificaci\'on matem\'atica}

\begin{tcolorbox}[math, title=Par\'ametros de este paso]
\textbf{Cero.} El horizonte (2h, 4h, 6h) es arbitrario y no sensible ---
si la se\~nal tiene estructura, la bimodalidad aparece en cualquier
horizonte razonable. Si solo aparece con un horizonte espec\'ifico,
la estructura es fr\'agil y no vale la pena.
\end{tcolorbox}

Necesit\'as al menos 50 trades para que los tests del Paso~2 tengan
potencia estad\'istica suficiente. Con la frecuencia t\'ipica de se\~nales
en altcoins perpetuos, eso es 2--4 semanas.


%=================================================================
\section{Paso 2: ?`Hay Estaci\'on? --- Test de Bimodalidad}
\label{sec:paso2}
%=================================================================

\begin{tcolorbox}[radio, title=La radio: ?`hay estaci\'on?]
Despu\'es de escuchar 50 frecuencias, revis\'as tus anotaciones.
En la mayor\'ia de las frecuencias, todo lo que anotaste es ruido
uniforme --- no hay patr\'on. Pero en algunas frecuencias, las
anotaciones muestran dos tipos de sonido claramente distintos:
a veces escuchaste algo que parec\'ia m\'usica (se\~nal) y a veces
est\'atica (ruido). Esa frecuencia tiene una estaci\'on.

Si todo lo que anotaste en una frecuencia es un continuo uniforme
sin dos grupos distinguibles, no hay estaci\'on --- es solo ruido
con variaciones aleatorias. No importa cu\'anto tiempo escuches,
nunca vas a decodificar un mensaje porque no hay mensaje.
\end{tcolorbox}

\subsection{El test}

Con los 50+ trades del Paso~1, construimos la distribuci\'on de MFE
(M\'axima Excursi\'on Favorable) y preguntamos: ?`esta distribuci\'on
tiene una joroba o dos?

Formalmente, comparamos dos modelos:
\begin{align}
M_1&: \text{MFE} \sim \mathcal{N}(\mu, \sigma^2) \quad \text{(una campana)} \\
M_2&: \text{MFE} \sim \pi_1 \mathcal{N}(\mu_1, \sigma_1^2) + \pi_2 \mathcal{N}(\mu_2, \sigma_2^2) \quad \text{(dos campanas)}
\end{align}

El BIC (Bayesian Information Criterion) de cada modelo es:
\begin{equation}
\text{BIC} = -2 \ln \hat{L} + k \ln n
\end{equation}
donde $\hat{L}$ es la verosimilitud maximizada, $k$ es el n\'umero de
par\'ametros ($k=2$ para $M_1$, $k=5$ para $M_2$), y $n$ es el n\'umero
de trades. El BIC penaliza la complejidad: $M_2$ tiene m\'as par\'ametros,
as\'i que necesita ajustar sustancialmente mejor para ``ganar''.

\begin{definition}[$\Delta$BIC]
\begin{equation}
\Delta\text{BIC} = \text{BIC}(M_1) - \text{BIC}(M_2)
\end{equation}
Si $\Delta\text{BIC} > 0$, dos campanas son mejor. Si $< 0$, una alcanza.
\end{definition}

\subsection{Escala de interpretaci\'on (Kass y Raftery, 1995)}

\begin{table}[H]
\centering
\begin{tabular}{ccl}
\toprule
$\Delta$\textbf{BIC} & \textbf{Evidencia} & \textbf{Acci\'on} \\
\midrule
$< 2$ & Insignificante & Descartar la se\~nal \\
$2$--$6$ & Positiva & Dudosa, seguir acumulando datos \\
$6$--$10$ & Fuerte & Prometedor, avanzar con cautela \\
$> 10$ & Muy fuerte & \textbf{Hay estaci\'on --- avanzar al Paso~3} \\
\bottomrule
\end{tabular}
\end{table}

\subsection{El criterio de falsificaci\'on}

\begin{tcolorbox}[warning, title=Regla inquebrantable]
Si $\Delta$BIC $< 2$ despu\'es de 50+ trades, \textbf{descart\'a la
se\~nal}. No hay discusi\'on, no hay ``probemos con otros par\'ametros'',
no hay optimizaci\'on. Si no hay dos campanas, no hay estructura
explotable. No hay estaci\'on en esa frecuencia. Mov\'e el dial.
\end{tcolorbox}

\subsection{Por qu\'e esto funciona: la cadena causal}

\begin{tcolorbox}[math, title=Teorema: bimodalidad es condici\'on necesaria]
\begin{theorem}
Sin bimodalidad en MFE, el estimador Shannon-VOI colapsa a
clasificaci\'on aleatoria.
\end{theorem}

\textbf{Demostraci\'on emp\'irica} (5 escenarios simulados, $N=83$):

\medskip
\begin{tabular}{lccc}
\toprule
\textbf{Escenario} & $\Delta$\textbf{BIC} & \textbf{I(MFE; O)} & \textbf{Separaci\'on} \\
\midrule
Real (bimodal) & $+13.9$ & 0.353 bits & $+44$ pp \\
Shuffled & $+13.9$ & 0.018 bits & $+6$ pp \\
Unimodal pura & $-5.9$ & 0.014 bits & $+16$ pp \\
Unimodal + correlaci\'on & $-7.2$ & 0.050 bits & $-1$ pp \\
RSI-like & $+10.5$ & 0.080 bits & $-6$ pp \\
\bottomrule
\end{tabular}

\medskip
Sin bimodalidad, la informaci\'on del canal cae de 0.353 a 0.014
bits (25$\times$ menor). La superficie $\pw$ se aplana y Shannon no
puede clasificar.
\end{tcolorbox}

\subsection{Par\'ametros de este paso}

\textbf{Cero.} El $\Delta$BIC es un estad\'istico --- sale de los datos
directamente. El umbral de 10 viene de la teor\'ia estad\'istica
bayesiana (Kass-Raftery), no de una optimizaci\'on. No hay nada que
calibrar.


%=================================================================
\section{Paso 3: Mapear la Se\~nal --- Superficie P(win $|$ MFE$_t$, $t$)}
\label{sec:paso3}
%=================================================================

\begin{tcolorbox}[radio, title=La radio: mapear la se\~nal]
Encontraste una estaci\'on. Ahora necesit\'as entender c\'omo
transmite. Escuch\'as muchas transmisiones y anot\'as: ``despu\'es
de 30 segundos, si la se\~nal tiene tal intensidad, el mensaje
suele ser X''. ``Despu\'es de 2 minutos, si la intensidad es
baja, el mensaje suele ser Y''. Vas construyendo un mapa: para
cada combinaci\'on de (tiempo transcurrido, intensidad de se\~nal),
sab\'es qu\'e mensaje es m\'as probable.

El mapa es una tabla 2D: filas = tiempo, columnas = intensidad,
valor en cada celda = probabilidad de que el mensaje sea ``W''.
\end{tcolorbox}

\subsection{Construcci\'on}

Usamos los path data del Paso~1 para construir la superficie
$P(\text{win} \mid \text{MFE}_t, t)$. Para cada trade, recorremos
su trayectoria y en cada punto $(t_i, \text{MFE}_i)$ registramos
si el trade termin\'o como winner o loser.

La superficie se construye como un histograma 2D suavizado:

\begin{enumerate}[leftmargin=2em]
\item Discretizar el espacio $(t, \text{MFE})$ en celdas.
\item Para cada celda, contar cu\'antos winners y losers pasaron por ah\'i.
\item $\pw$ en cada celda = winners / (winners + losers).
\item Suavizar con kernel gaussiano (bandwidth de Scott).
\end{enumerate}

\subsection{Qu\'e revela la superficie}

Si la se\~nal tiene estructura (Paso~2 positivo), la superficie tiene
\textbf{gradiente fuerte}: las celdas con MFE alto tienen $\pw$ alto
(70--90\%) y las celdas con MFE bajo tienen $\pw$ bajo (30--40\%).
Las l\'ineas de iso-probabilidad son curvas diagonales en el espacio
$(t, \text{MFE})$.

\begin{table}[H]
\centering
\caption{Ejemplo real: $\pw$ en coordenadas seleccionadas.}
\begin{tabular}{r|ccccc}
\toprule
& \multicolumn{5}{c}{\textbf{MFE al tiempo $t$}} \\
$t$ (min) & 0.0\% & 0.4\% & 0.8\% & 1.2\% & 2.0\% \\
\midrule
0.5 & 49\% & 53\% & 57\% & 62\% & 75\% \\
2.0 & 47\% & 52\% & 56\% & 62\% & 74\% \\
5.0 & 43\% & 49\% & 54\% & 59\% & 70\% \\
10.0 & 41\% & 45\% & 48\% & 54\% & 66\% \\
20.0 & 36\% & 41\% & 46\% & 50\% & 55\% \\
\bottomrule
\end{tabular}
\end{table}

Si la se\~nal \textbf{no} tiene estructura (Paso~2 negativo), la
superficie es \textbf{plana}: todas las celdas tienen $\pw \approx$
win rate base ($\sim$50\%). No hay gradiente, no hay curvas, no hay
nada que explotar. Shannon mira esta tabla y ve: ``todo es igual,
no puedo distinguir nada''.

\subsection{Justificaci\'on matem\'atica: informaci\'on mutua}

La informaci\'on que la superficie captura se mide con la
informaci\'on mutua de Shannon:

\begin{equation}
I(\text{MFE}_t; O) = H(O) - H(O \mid \text{MFE}_t)
\end{equation}

donde $H(O)$ es la entrop\'ia del outcome a priori:
\begin{equation}
H(O) = -\pi_W \log_2 \pi_W - (1-\pi_W) \log_2(1-\pi_W)
\end{equation}

y $H(O \mid \text{MFE}_t)$ es la entrop\'ia residual despu\'es de
observar el MFE al tiempo $t$:
\begin{equation}
H(O \mid \text{MFE}_t) = \sum_{\text{celda } c} P(c) \cdot H_b\!\left(\pw_c\right)
\end{equation}

\begin{tcolorbox}[result, title=Resultado emp\'irico: ratio 13:1]
En nuestros datos ($N=83$ trades):
\begin{itemize}[leftmargin=1.5em]
\item La se\~nal de entrada aporta: $I_{\text{se\~nal}} = 0.013$ bits
\item La trayectoria del precio aporta: $I_{\text{trayectoria}} = 0.167$ bits
\item \textbf{Ratio: 13:1 a favor de la observaci\'on post-entrada}
\end{itemize}

La gran mayor\'ia de la informaci\'on relevante no est\'a disponible
al momento de entrar --- se revela progresivamente a trav\'es del canal.
Esto justifica formalmente por qu\'e la gesti\'on de salida importa
13 veces m\'as que la se\~nal de entrada.
\end{tcolorbox}

\subsection{Par\'ametros de este paso}

\textbf{Cero.} La superficie es un conteo por celda (histograma 2D),
no un modelo con par\'ametros. El \'unico choice es el bandwidth del
kernel de suavizado, y el default de Scott funciona.


%=================================================================
\section{Paso 4: Decodificar el Mensaje --- Sistema de Salida Shannon-VOI}
\label{sec:paso4}
%=================================================================

\begin{tcolorbox}[radio, title=La radio: decodificar el mensaje]
Ten\'es el mapa de la estaci\'on (la superficie). Ahora sintoniz\'as
una nueva transmisi\'on (abr\'is un trade) y empez\'as a escuchar.

En cada momento, mir\'as el mapa y dec\'is: ``llevo 3 minutos
escuchando, la intensidad de la se\~nal es media-baja... seg\'un
el mapa, esto tiene 45\% de probabilidad de ser un mensaje~W''.

Tres situaciones posibles:

\textbf{1. La se\~nal es d\'ebil y dej\'o de cambiar.} Llevás 10
minutos, la intensidad es baja, y no cambia. Seg\'un el mapa,
$\pw < 50\%$ y la tasa de informaci\'on nueva es cero. La estaci\'on
ya dijo todo lo que ten\'ia que decir, y el mensaje es ``L''. Dej\'as
de escuchar --- no tiene sentido seguir.

\textbf{2. La se\~nal es fuerte y estable.} Llevás 5 minutos, la
intensidad es alta, $\pw > 70\%$. El mensaje es ``W'' con alta
confianza. Ahora no te preocupa \emph{si} es winner sino
\emph{cu\'anto} dura. Sub\'is el volumen (trailing activo) y
esper\'as a que la se\~nal se debilite naturalmente.

\textbf{3. La se\~nal todav\'ia est\'a cambiando.} Llevás 2 minutos
y la intensidad sube y baja. La tasa de informaci\'on nueva es alta
--- la estaci\'on todav\'ia est\'a transmitiendo datos \'utiles. No
tom\'as ninguna decisi\'on --- segu\'is escuchando porque cada segundo
adicional reduce tu incertidumbre.
\end{tcolorbox}

\subsection{Las tres zonas de decisi\'on}

\begin{tcolorbox}[math, title=Criterio \'optimo de parada (Wald-Shannon)]
En cada momento $t$ del trade, el sistema calcula:
\begin{itemize}[leftmargin=1.5em]
\item $\pw(t)$: probabilidad de winner, interpolada de la superficie
\item $dI/dt$: tasa de informaci\'on nueva (bits por minuto)
\item $\E[\text{PnL} \mid \text{continuar}] = \pw \cdot \E[W] + (1-\pw) \cdot \E[L]$
\end{itemize}

\medskip
\textbf{Zona 1 --- Aprendizaje} ($dI/dt > \varepsilon$, cualquier $\pw$):\\
HOLD. El canal todav\'ia transmite informaci\'on \'util. Decidir ahora
desperdicia informaci\'on que est\'a llegando gratis.

\medskip
\textbf{Zona 2 --- Decisi\'on negativa} ($dI/dt \leq \varepsilon$, $\pw < 0.50$):\\
EXIT. La informaci\'on se agot\'o y dice ``loser''. Esperar m\'as no
cambiar\'a el veredicto porque $dI/dt \approx 0$.

\medskip
\textbf{Zona 3 --- Decisi\'on positiva} ($dI/dt \leq \varepsilon$, $\pw > 0.70$):\\
TRAIL. La informaci\'on se agot\'o y dice ``winner''. Activar trailing
con callback $= f(\pw)$ para capturar el recorrido.
\end{tcolorbox}

\subsection{Por qu\'e $\pw = 0.50$ no es un par\'ametro}

El umbral de EXIT en $\pw = 0.50$ no es una elecci\'on --- es una
consecuencia matem\'atica. El valor esperado de continuar es:
\begin{equation}
\E[\text{PnL} \mid \text{continuar}] = \pw \cdot \E[\text{PnL} \mid W] + (1-\pw) \cdot \E[\text{PnL} \mid L]
\end{equation}

Con $\E[W] = +2.54\%$ y $\E[L] = -3.08\%$:
\begin{equation}
\E = 0 \quad\Longleftrightarrow\quad \pw = \frac{|\E[L]|}{|\E[W]| + |\E[L]|} = \frac{3.08}{2.54 + 3.08} = 0.548
\end{equation}

El punto exacto donde $\E = 0$ es 0.548, cercano a 0.50. Usar 0.50
es ligeramente conservador (cierra un poco antes del punto de EV
cero), lo cual es correcto porque incluye un margen por costos de
transacci\'on y slippage.

\subsection{El trailing adaptativo}

Cuando $\pw > 0.70$, el sistema activa trailing con callback que
depende de $\pw$:
\begin{equation}
\text{callback}(\pw) = 0.30 + (\pw - 0.70) \times 0.833
\end{equation}

Esto da: $\pw = 0.70 \to 30\%$, $\pw = 0.80 \to 38\%$,
$\pw = 0.90 \to 47\%$, $\pw = 0.95 \to 51\%$.

Mayor confianza $\to$ callback m\'as suelto (dejar correr).
Menor confianza $\to$ callback m\'as apretado (proteger).

\subsection{Justificaci\'on te\'orica: Data Processing Inequality}

?`Por qu\'e este sistema de un solo par\'ametro supera a uno de 14?

\begin{tcolorbox}[math, title=Data Processing Inequality (Shannon 1948)]
Para cualquier cadena de Markov $X \to Y \to Z$:
\begin{equation}
I(X; Z) \leq I(X; Y)
\end{equation}
Ninguna transformaci\'on de los datos puede \emph{crear} informaci\'on
que no estaba en los datos originales. Solo puede destruirla.
\end{tcolorbox}

El AEPS toma la trayectoria continua $Y$ y la reduce a comparaciones
binarias $Z$: ``?`MFE $> 1.14\%$?'', ``?`MAE $> 5.3\%$?''. Cada
threshold es una compresi\'on lossy que destruye informaci\'on.

El estimador Shannon opera directamente sobre $Y$ --- usa la
distribuci\'on continua $\pw$ en vez de reducirla a ``cruz\'o/no
cruz\'o''. Por el DPI, captura m\'as informaci\'on y clasifica mejor.

\begin{tcolorbox}[result, title=Resultado emp\'irico]
\begin{center}
\begin{tabular}{lcc}
\toprule
& \textbf{AEPS (14 par\'ametros)} & \textbf{Shannon (1 par\'ametro)} \\
\midrule
PnL acumulado & $+8.6\%$ & $+35.8\%$ \\
Mejora & --- & $4.2\times$ \\
\bottomrule
\end{tabular}
\end{center}
La mejora de $4.2\times$ es consistente con el DPI: preservar
informaci\'on en vez de destruirla con thresholds produce mejor
clasificaci\'on.
\end{tcolorbox}

\subsection{Par\'ametros de este paso}

\textbf{Uno:} trail\_pw $= 0.70$ (cu\'ando activar trailing). Es robusto:
moverlo a 0.65 o 0.75 cambia cu\'ando se activa el trailing por unos
minutos, y el callback adaptativo $f(\pw)$ absorbe la diferencia.

Adem\'as hay dos mecanismos que sobreviven independientes de Shannon:
\begin{itemize}[leftmargin=1.5em]
\item \textbf{Hard SL} $= 8\%$: protecci\'on de capital contra cisnes negros.
  No es clasificaci\'on --- es seguro.
\item \textbf{Pump capture}: se\~nal de mercado ($d$OI $< 0$, taker ratio
  $< 45\%$) que detecta exhaustion. Shannon no ve esta informaci\'on
  porque solo mira MFE/PnL/MAE.
\end{itemize}


%=================================================================
\section{Paso 5: Monitorear la Estaci\'on --- Detecci\'on de Erosi\'on}
\label{sec:paso5}
%=================================================================

\begin{tcolorbox}[radio, title=La radio: la estaci\'on se debilita]
Ven\'is escuchando la estaci\'on hace semanas y funciona bien.
Pero un d\'ia not\'as que la se\~nal es m\'as d\'ebil, m\'as ruidosa,
m\'as dif\'icil de decodificar. Los mensajes que antes eran claros
ahora son ambiguos. La estaci\'on no desapareci\'o de golpe --- se
fue debilitando gradualmente.

Necesit\'as un medidor que te diga: ``la potencia de la estaci\'on
cay\'o del nivel X al nivel Y''. Cuando baja de cierto nivel, dej\'as
de escuchar y busc\'as otra.
\end{tcolorbox}

\subsection{Monitoreo por $\Delta$BIC rolling}

Cada 50 trades, recalcul\'as el $\Delta$BIC sobre los \'ultimos
50--100 trades. Si $\Delta$BIC cae por debajo de 6, la bimodalidad
se est\'a comprimiendo --- el primer eslab\'on de la cadena se est\'a
debilitando. Si cae por debajo de 2, la bimodalidad desapareci\'o
y el edge se perdi\'o.

\subsection{Monitoreo online: filtro Beta-Binomial}

Para detecci\'on trade a trade (sin esperar a 50), implementamos un
filtro bayesiano que trackea $\pi_W$ (win rate) en tiempo real:

\begin{equation}
\text{Si win:}\quad \alpha_{t+1} = \lambda\alpha_t + 1, \quad \beta_{t+1} = \lambda\beta_t
\end{equation}
\begin{equation}
\text{Si loss:}\quad \alpha_{t+1} = \lambda\alpha_t, \quad \beta_{t+1} = \lambda\beta_t + 1
\end{equation}

donde $\lambda = 0.97$ es un factor de descuento (half-life $\approx$
23 trades). La estimaci\'on de $\pi_W$ es:
\begin{equation}
\hat{\pi}_W = \frac{\alpha}{\alpha + \beta}
\end{equation}

\begin{tcolorbox}[warning, title=Alertas]
\begin{itemize}[leftmargin=1.5em]
\item $\hat{\pi}_W < 0.45$ durante 20+ trades $\to$ el edge puede
  estar erosion\'andose. Investigar.
\item $\Delta$BIC $< 6$ en ventana rolling $\to$ la bimodalidad se
  comprime. Regenerar superficie y evaluar.
\item $\Delta$BIC $< 2$ $\to$ \textbf{pausar}. No hay estaci\'on.
\end{itemize}
\end{tcolorbox}

\subsection{Actualizaci\'on incremental de la superficie}

La superficie se actualiza despu\'es de cada trade cerrado, no en
batch cada 50 trades. Cada celda $(t, \text{MFE})$ mantiene sus
propios contadores Beta$(\alpha_c, \beta_c)$:

\begin{equation}
\text{Para cada punto del path del trade cerrado:} \quad
\begin{cases}
\alpha_c \leftarrow \lambda\alpha_c + 1 & \text{si winner} \\
\beta_c \leftarrow \lambda\beta_c + 1 & \text{si loser}
\end{cases}
\end{equation}

\begin{equation}
\pw_c = \frac{\alpha_c}{\alpha_c + \beta_c}
\end{equation}

La superficie nunca est\'a ``vieja'' --- se actualiza en cada trade.
Con $\lambda = 0.97$, trades de hace 50 operaciones tienen peso
$0.97^{50} = 22\%$ del peso de un trade reciente. La superficie se
adapta gradualmente sin discontinuidades.

\subsection{Par\'ametros de este paso}

\textbf{Cero operativos.} Los umbrales de alerta ($\Delta$BIC $< 6$,
$\pi_W < 0.45$) vienen de la teor\'ia estad\'istica, no de
optimizaci\'on. El decay $\lambda = 0.97$ da un half-life de
$\sim$23 trades, que es un balance razonable entre reactividad y
estabilidad --- pero no es sensible a $\pm$0.02.


%=================================================================
\section{Resumen: De 14 Par\'ametros a 1}
\label{sec:resumen}
%=================================================================

\begin{tcolorbox}[key, title=El framework completo]
\begin{center}
\renewcommand{\arraystretch}{1.4}
\begin{tabular}{clcc}
\toprule
\textbf{Paso} & \textbf{Acci\'on} & \textbf{Radio} & \textbf{Par\'ametros} \\
\midrule
1 & Se\~nal $\to$ dry mode horizonte fijo & Girar el dial & 0 \\
2 & $\Delta$BIC $> 10$? & ?`Hay estaci\'on? & 0 \\
3 & Superficie $\pw$ & Mapear la se\~nal & 0 \\
4 & Shannon EXIT / TRAIL / HOLD & Decodificar & \textbf{1} \\
5 & $\Delta$BIC rolling + $\hat{\pi}_W$ & Monitorear potencia & 0 \\
\midrule
& \textbf{Total} & & \textbf{1} \\
\bottomrule
\end{tabular}
\end{center}
\end{tcolorbox}

Comparaci\'on con el sistema anterior:

\begin{center}
\renewcommand{\arraystretch}{1.3}
\begin{tabular}{lcc}
\toprule
& \textbf{AEPS} & \textbf{Shannon} \\
\midrule
Par\'ametros calibrados & 14 & 1 \\
Necesita backtesting & S\'i & No \\
Necesita optimizaci\'on & S\'i & No \\
Criterio de falsificaci\'on & No & S\'i ($\Delta$BIC) \\
Criterio de muerte & No & S\'i ($\Delta$BIC $< 2$) \\
Agnóstico a la se\~nal & No & S\'i \\
Agnóstico al mercado & No & S\'i \\
Resultado ($N=83$) & $+8.6\%$ & $+35.8\%$ \\
\bottomrule
\end{tabular}
\end{center}


%=================================================================
\section{?`Por Qu\'e Funciona con Menos?}
\label{sec:porquemenos}
%=================================================================

\begin{tcolorbox}[radio, title=La radio: por qu\'e un receptor simple es mejor]
Un receptor de radio sofisticado con 14 perillas (graves, agudos,
balance, filtro de ruido, AGC, squelch, etc.) puede sonar bien si
calibr\'as cada perilla perfectamente. Pero si una perilla est\'a
mal, distorsiona todo. Y cuando la estaci\'on cambia de formato,
ten\'es que recalibrar las 14.

Un receptor simple que solo tiene ``volumen'' funciona peor en el
caso \'optimo, pero no se rompe nunca. Y si le das acceso directo
a la se\~nal en vez de pasarla por 14 filtros, captura m\'as de la
transmisi\'on original.

Eso es exactamente el Data Processing Inequality: cada perilla es
un filtro que destruye informaci\'on. Menos filtros = m\'as se\~nal.
\end{tcolorbox}

Los 14 par\'ametros del AEPS eran aproximaciones discretas de una
funci\'on continua (la superficie $\pw$). Cada par\'ametro capturaba
un pedazo de esa funci\'on:
\begin{itemize}[leftmargin=1.5em]
\item El abort capturaba la zona $\pw < 0.40$
\item El partial TP capturaba la zona $\pw \approx 0.55$--$0.65$
\item El trailing capturaba la zona $\pw > 0.70$
\item El zombie kill capturaba la zona $\pw < 0.50$ con $t > t_{\text{abort}}$
\end{itemize}

Al discretizar, perd\'ias informaci\'on (DPI) y necesitabas m\'as
par\'ametros para compensar la p\'erdida. Con la funci\'on continua,
no necesit\'as compensar nada.

Es la diferencia entre aproximar una curva con 14 segmentos rectos
versus tener la curva misma. Los segmentos necesitan 28 coordenadas
(14 posiciones $\times$ 2 extremos). La curva necesita 0 par\'ametros
--- ES la curva, directamente de los datos.


%=================================================================
\section{Conexi\'on con la F\'isica: Weibull y Shannon Dicen lo Mismo}
\label{sec:weibull_shannon}
%=================================================================

\begin{tcolorbox}[radio, title=La radio: por qu\'e el mensaje llega al principio]
Cuando sintoniz\'as una estaci\'on, los primeros segundos son los
m\'as informativos. Al principio no sab\'es nada --- podr\'ia ser
m\'usica, noticias, o est\'atica. Pero r\'apidamente capt\'as el tono:
``esto es m\'usica'' o ``esto es ruido''. Despu\'es de 2 minutos, ya
sab\'es. Despu\'es de 10 minutos, escuchar m\'as no agrega nada.

La raz\'on es f\'isica: la estaci\'on transmite a una potencia que
decae con el tiempo. Los primeros bits de informaci\'on llegan con
se\~nal fuerte. Los siguientes, m\'as d\'ebiles. Eventualmente la
se\~nal se pierde en el ruido y no hay m\'as informaci\'on por extraer.
\end{tcolorbox}

La supervivencia Weibull ($k = 0.441 < 1$) y la tasa de informaci\'on
de Shannon ($dI/dt$) describen el mismo fen\'omeno con lenguajes
diferentes:

\begin{itemize}[leftmargin=1.5em]
\item \textbf{Weibull dice:} la probabilidad de escape (la transici\'on
  informativa) decae como $h(t) \propto t^{k-1} = t^{-0.56}$.
\item \textbf{Shannon dice:} la tasa de informaci\'on nueva decae de
  0.05 bits/min a los 2 min hasta 0.001 bits/min a los 30 min.
\end{itemize}

Son dos caras de la misma moneda, y esa moneda es la ecuaci\'on de
Langevin con su discriminante de damping.

\begin{tcolorbox}[math, title=Conexi\'on formal]
El estado latente del trade $Z \in \{W, L\}$ se revela a trav\'es de
un evento de escape (cruce de barrera en el potencial de Langevin).
La probabilidad de este evento en el intervalo $[t, t+dt]$ es la
hazard $h(t)$.

La informaci\'on ganada por observar si el escape ocurri\'o o no es:
\begin{equation}
dI \propto h(t) \cdot dt \cdot D_{\text{KL}}(f_W \| f_L)
\end{equation}

Con $h(t) = \frac{k}{\lambda}(t/\lambda)^{k-1}$ y $k < 1$, la hazard
(y por lo tanto la tasa de informaci\'on) decae monotónicamente.
La informaci\'on llega r\'apido al principio y se agota.

Esto justifica el zombie kill (no esperar una transmisi\'on que ya se
agot\'o) y el trailing temprano (actuar cuando la informaci\'on dice
``winner'' con confianza, no esperar a una confirmaci\'on que no va a
llegar).
\end{tcolorbox}


%=================================================================
\section{Aplicabilidad: ?`D\'onde Funciona Este Framework?}
\label{sec:aplicabilidad}
%=================================================================

El framework es agn\'ostico a la se\~nal y al mercado, pero requiere
tres condiciones:

\begin{enumerate}[leftmargin=2em]
\item \textbf{Path data observable.} Se necesita la trayectoria del
  precio tick a tick, no solo el resultado final. Mercados con datos
  de alta frecuencia accesibles: cripto perpetuos, forex, futuros
  regulados. No funciona en mercados donde solo ves el precio de
  cierre diario.

\item \textbf{Microestructura mec\'anica.} La bimodalidad es m\'as
  probable en mercados con fuerzas mec\'anicas (funding, liquidaciones,
  margin calls) que crean bifurcaciones binarias. Mercados candidatos:
  perpetuos en dYdX/GMX, futuros en CME, opciones con gamma squeeze.

\item \textbf{Frecuencia suficiente.} Se necesitan 50+ trades para
  el test de bimodalidad. Se\~nales que disparan 1 trade por mes no
  acumulan datos suficientes en tiempo razonable.
\end{enumerate}


%=================================================================
\section{Conclusi\'on: El Producto es la Metodolog\'ia}
\label{sec:conclusion}
%=================================================================

\begin{tcolorbox}[key, title=Lo que se vende no es un bot]
El producto no es un bot que gana plata shorteando altcoins. Es un
\textbf{detector de estaciones} con un \textbf{decodificador autom\'atico}:

\begin{itemize}[leftmargin=1.5em]
\item El detector ($\Delta$BIC) funciona en cualquier banda (mercado).
\item El decodificador (Shannon-VOI) funciona con cualquier estaci\'on
  (se\~nal bimodal).
\item El monitor ($\hat{\pi}_W$ + $\Delta$BIC rolling) detecta cuando
  la estaci\'on se debilita.
\end{itemize}

El framework es agn\'ostico a la frecuencia --- funciona sin importar
qu\'e estaci\'on encontr\'es, siempre que haya una estaci\'on.
\end{tcolorbox}

Propiedades del framework:
\begin{itemize}[leftmargin=1.5em]
\item \textbf{Falsificable:} $\Delta$BIC $< 2$ $\to$ no hay edge. No hay ambig\"uedad.
\item \textbf{Minimal:} 1 par\'ametro operativo. Auditable por cualquiera.
\item \textbf{Sin backtesting:} la superficie sale de datos en vivo, no de
  historia optimizada.
\item \textbf{Con criterio de muerte:} el sistema sabe cu\'ando dejar de
  operar.
\item \textbf{Te\'oricamente fundado:} Langevin $\to$ Weibull $\to$
  Shannon $\to$ VOI. Cada eslab\'on con validaci\'on emp\'irica.
\end{itemize}


%=================================================================
\begin{thebibliography}{10}
%=================================================================

\bibitem{shannon1948} C.~E.~Shannon, ``A Mathematical Theory of Communication,'' \textit{Bell Syst.\ Tech.\ J.}, 27(3):379--423, 1948.

\bibitem{kass1995} R.~E.~Kass y A.~E.~Raftery, ``Bayes Factors,'' \textit{JASA}, 90(430):773--795, 1995.

\bibitem{wald1945} A.~Wald, ``Sequential Tests of Statistical Hypotheses,'' \textit{Ann.\ Math.\ Stat.}, 16(2):117--186, 1945.

\bibitem{cover2006} T.~M.~Cover y J.~A.~Thomas, \textit{Elements of Information Theory}, 2da ed., Wiley, 2006.

\bibitem{howard1966} R.~A.~Howard, ``Information Value Theory,'' \textit{IEEE Trans.\ Syst.\ Sci.\ Cybern.}, 2(1):22--26, 1966.

\bibitem{kaplan1958} E.~L.~Kaplan y P.~Meier, ``Nonparametric Estimation from Incomplete Observations,'' \textit{JASA}, 53(282):457--481, 1958.

\bibitem{kramers1940} H.~A.~Kramers, ``Brownian Motion in a Field of Force,'' \textit{Physica}, 7(4):284--304, 1940.

\end{thebibliography}

\end{document}
